\section{When Not to Start a Company}

\subsection{The Wrong Reasons to Found}

It would be irresponsible to end a chapter on entrepreneurship without considering the cases in which entrepreneurship is a poor immediate choice.

You probably should not start a full-time, high-burn company merely because you dislike your current manager. Changing jobs may solve the problem with far less cost.

You probably should not borrow heavily to launch a business whose customer demand has not been tested, especially when the borrowing can compromise essential household obligations. A business experiment is not worth destroying the ability to make future experiments.

You probably should not abandon a rare opportunity to learn from extraordinary colleagues merely because the founder title sounds more impressive. Some employment positions offer an unusually high return in skills, credibility, networks, and savings. Those assets may become the foundation of a much stronger business later.

You probably should not accept venture capital because investors are willing to offer it. Their financing has a logic: it often assumes rapid growth, substantial scale, and a return large enough to compensate for many failed investments. If you seek a steady cash-generating company or a mission-driven institution with slow capital formation, venture-style incentives can create damaging pressure.

And you probably should not begin with an enterprise whose daily operation you fundamentally dislike. The famous founder may spend time thinking about strategy. The actual early founder spends time responding to confused customers, correcting mistakes, collecting payments, explaining the product repeatedly, writing mundane documents, and fixing systems at inconvenient hours. The character of the work matters more than its public mythology.

These are not prohibitions. A person can knowingly accept extraordinary risk for a deeply valued goal. But that risk should be chosen deliberately, not stumbled into through imitation or social pressure.

\subsection{The Disciplined Alternative: Constrained Experimentation}

The rational alternative to going all-in immediately is not necessarily doing nothing. It may be a deliberately constrained experiment. Build a prototype. Sell a small service. Join an organization close to the target industry. Publish serious work. Save money. Develop relationships. Find a collaborator. Test whether people will exchange resources for the value you propose. Each action can convert uncertainty into information and accumulate an asset even if it does not produce a successful company.

There are circumstances where waiting is itself costly (an opening in a fast-changing technology market may close), but the cost of waiting should be evaluated rather than assumed. \textbf{Start immediately and never take risks are equally poor universal rules.}
